\documentclass[10pt,a4paper]{amsart}
\usepackage[margin=19mm]{geometry}
\usepackage{amsmath,amssymb,mathtools}
\usepackage[scr=rsfs,cal=euler]{mathalfa}
\usepackage{xfrac}
\usepackage[unicode,hidelinks,bookmarks=false]{hyperref}
\newcommand{\ie}{i.e.\ }
\newcommand{\cf}{cf.\ }
\newcommand{\resp}{respectively\ }
\newcommand{\Subsection}{\subsection}
\providecommand{\nobreakdash}{\nobreak\hbox{-}}
\setlength{\emergencystretch}{2em}
\multlinegap0pt
\usepackage{bm}
\usepackage{bbm}
\usepackage{dsfont}
\usepackage{xspace}
\usepackage{cancel}
\usepackage{mathtools}
\usepackage{cleveref}
\usepackage{amsthm}
\makeatletter
\@ifundefined{theorem}{\newtheorem{theorem}{Theorem}}{}
\makeatother
\title[Lions and Contamination: Trees and General Graphs]{Lions and Contamination: Trees and General Graphs}
\author{Dohoon Kim, Eungyu Woo, Donghoon Shin}
\date{Source: arXiv:2604.18949v1 (CC BY 4.0)}
\begin{document}
\maketitle

\noindent\emph{Source and licence.} Lions and Contamination: Trees and General Graphs. Version arXiv:2604.18949v1 (CC BY 4.0), distributed under the Creative Commons Attribution 4.0 International licence, as recorded in the arXiv licence metadata for this exact version and shown on the abstract page https://arxiv.org/abs/2604.18949v1. Full rights notice: licenses/arxiv-2604.18949-CC-BY-4.0.txt.\par\smallskip

\noindent\textit{Editorial scope.} Counted unit: the Theorem whose source label is \texttt{theorem:subgraph\_counterexample} in the article (source0.tex lines 299--301, source label \texttt{theorem:subgraph\_counterexample}), delivered together with the article's own complete original proof of it (source0.tex lines 303--322, one \texttt{proof} environment of 20 lines). No mathematics is written, repaired or re-derived: statement and proof are the article's own text line for line. Required context retained uncounted: theorem:pathwidth\_theorem (lines 271--273). Every definition, display and label of those lines is retained unchanged. Omitted from this package and not redistributed: 1--270, 274--298, 302--302, 323--555. Nothing inside a retained excerpt is omitted or altered: every retained body is delivered exactly as the article's lines are written, so no omission receipt of any kind accompanies this package. Citations: the retained excerpts carry no citation command at all, so no bibliography entry of the article is transcribed and no reference is redistributed. Reference adaptation: none was needed and none was made; every reference inside the delivered unit resolves inside it, so no unresolved reference or citation is produced. Printed slips kept literal: no printed word, symbol, hypothesis or number of the counted statement or of its proof was altered. Layout and class: the article's own class is \texttt{lipics-v2021} with packages including ; the wrapper is the lane's pinned \texttt{amsart} editorial wrapper at 10pt on A4, which restates the theorem-like environments the retained text needs and adds the article's own macro definitions that the retained text needs (none), with extra wrapper packages (bm, bbm, dsfont, xspace, cancel, mathtools, cleveref). The delivered numbering is the unit's own. Assets: no retained span contains a figure environment, a graphics-inclusion command, a PDF-page inclusion, a table or a TikZ picture; no image file is copied into this package, the package declares no asset dependency and no image file is redistributed with it. Licence: version arXiv:2604.18949v1 of this article is distributed under the Creative Commons Attribution 4.0 International licence, as recorded in the arXiv licence metadata for this exact version and shown on the abstract page https://arxiv.org/abs/2604.18949v1. A full rights notice is delivered at licenses/arxiv-2604.18949-CC-BY-4.0.txt. Characters: every retained excerpt is pure ASCII, so the delivered engine, XeLaTeX with its default Latin Modern text font, reports no missing character.
\begin{theorem}\label{theorem:pathwidth_theorem}
For a tree $T$, $\operatorname{pw}(T) \leq \mathcal{L}(T) \leq \operatorname{pw}(T)+1$.
\end{theorem}

\begin{theorem}\label{theorem:subgraph_counterexample}
    The lion number is not monotonic under general subgraph inclusion.
\end{theorem}

\begin{proof}
To demonstrate this, we construct a family of trees $T_i$ and their supergraphs $G_i$ such that the supergraph requires fewer lions than its subgraph. Let $T_1$ be a tree consisting of a root and two leaves, and let $G_1$ be the graph formed by adding a universal vertex $u$ connected to all vertices in $T_1$. We define $\pi^{(1)}$ as the path that clears $G_1$ using exactly $3$ lions. In this strategy, one lion is placed on $u$ and the other two lions occupy the two leaves of $T_1$. In the subsequent time step, the two lions simultaneously slide to the root of $T_1$. This path $\pi^{(1)}$ requires $t_1 = 1$ time step, with the two lions arriving at the root of $T_1$ only at the final step.

We construct $T_{i+1}$ and $G_{i+1}$ inductively by assuming that a path $\pi^{(i)}$ clears $G_i$ using $3$ lions in $t_i$ time steps. We further assume that $\pi^{(i)}$ maintains one lion on the universal vertex $u$ throughout the process, while the other two lions arrive at the root of $T_i$ simultaneously at the last time step. The tree $T_{i+1}$ is constructed with a new root vertex $r_{i+1}$ connected to the roots of two disjoint copies of $T_i$, referred to as $T_L$ and $T_R$, via paths containing $t_i+5$ internal vertices. Consequently, the distance from $r_{i+1}$ to the root of $T_L$ is exactly $t_i+6$ edges. Let $G_{i+1}$ be the supergraph formed by adding a universal vertex $u$ connected to every vertex in $T_{i+1}$. A new path $\pi^{(i+1)}$ clears $G_{i+1}$ using $3$ lions according to the following steps:

\begin{enumerate}
    \item One lion remains at the universal vertex $u$ for the entire duration of the clearing process.
    \item The remaining two lions clear the left subtree $T_L$ by mimicking $\pi^{(i)}$.
    \item The two lions move up the path from the root of $T_L$ until they reach the child of $r_{i+1}$.
    \item Both lions jump to the universal vertex $u$, leaving the left path exposed to the contaminated root $r_{i+1}$.
    \item The lions jump from $u$ to the starting positions of $\pi^{(i)}$ within the right subtree $T_R$.
    \item The path $\pi^{(i)}$ is executed to clear $T_R$ in $t_i$ time steps, concluding with both lions at the root of $T_R$.
    \item One lion moves from the root of $T_R$ to the root of $T_L$ via $u$.
    \item The lions at the roots of $T_L$ and $T_R$ simultaneously slide up their respective paths, arriving at and clearing $r_{i+1}$ at the final time step of the process.
\end{enumerate}

It remains to verify that $T_L$ is not recontaminated during Steps 4 through 7. The total time elapsed between the lions vacating the left path in Step 4 and a lion returning to the root of $T_L$ in Step 7 is exactly $1 + 1 + t_i + 2 = t_i + 4$ time steps. Given that contamination spreads at a rate of one edge per time step, it advances at most $t_i+4$ edges down the left path from $r_{i+1}$. Since the total distance to the root of $T_L$ is $t_i+6$ edges, the contamination cannot reach $T_L$, ensuring it remains cleared. Furthermore, $\pi^{(i+1)}$ preserves the inductive condition where the lions arrive at the root $r_{i+1}$ only at the last time step.

Through this construction, we establish that $G_i$ can be cleared using only $3$ lions for an arbitrarily large $i$. Conversely, since $T_i$ contains a subdivision of a complete binary tree of height $i$, its pathwidth satisfies $\operatorname{pw}(T_i) \geq \lceil \frac{i}{2} \rceil$. According to \cref{theorem:pathwidth_theorem}, $\mathcal{L}(T_i) \geq \operatorname{pw}(T_i) \geq \lceil \frac{i}{2} \rceil$. For any $i > 6$, it follows that $\mathcal{L}(G_i) = 3 < \mathcal{L}(T_i)$. As $T_i$ is a subgraph of $G_i$, this confirms that the addition of edges or vertices to a graph can strictly decrease the required number of lions.
\end{proof}

\end{document}
